% Generated draft: requires mathematical and source review.
\documentclass[11pt]{article}
\usepackage[margin=25mm]{geometry}
\usepackage{fontspec}
\setmainfont{Latin Modern Roman}
\usepackage{amsmath,amssymb,mathtools,unicode-math}
\setmathfont{Latin Modern Math}
\usepackage{enumitem,needspace,xurl,hyperref}
\setlength{\parindent}{0pt}
\setlength{\parskip}{5pt}
\setlength{\emergencystretch}{3em}
\newcommand{\N}{\mathbb{N}}
\newcommand{\Z}{\mathbb{Z}}
\newcommand{\Q}{\mathbb{Q}}
\newcommand{\R}{\mathbb{R}}
\newcommand{\C}{\mathbb{C}}
\newcommand{\F}{\mathbb{F}}
\newcommand{\A}{\mathbb{A}}
\newcommand{\PP}{\mathbb{P}}
\newcommand{\E}{\mathbb{E}}
\newcommand{\eps}{\varepsilon}
\newcommand{\dd}{\,\mathrm{d}}
\begin{document}
\section*{UnsolvedMath ID: 10900015}
\textbf{Source problem code:} catalogue entry
\subsection{Definitions and mathematical statement}
A Kleinian group is a discrete subgroup \(\Gamma\leq\operatorname{PSL}_2(\C)\), acting by orientation-preserving isometries of hyperbolic \(3\)-space. In the usual closed-surface convention, a closed quasi-Fuchsian surface subgroup is the image of an injection
\[
 \rho:\pi_1(S_g)\hookrightarrow\Gamma,\qquad g\geq2,
\]
where \(S_g\) is a closed connected oriented genus-\(g\) surface and the image is quasi-Fuchsian of closed type, equivalently quasiconformally conjugate to a cocompact Fuchsian surface group. In particular, this convention does not replace the closed surface by a punctured surface.

\emph{Problem (Agol).} Characterize exactly the Kleinian groups \(\Gamma\) admitting such a subgroup:
\[
 \exists g\geq2\ \exists\rho:\pi_1(S_g)\hookrightarrow\Gamma
 \quad\rho(\pi_1(S_g))\text{ is closed quasi-Fuchsian}.
\]
The source identifies groups with parabolic subgroups as the difficult case. It reports that, without parabolics, the condition holds exactly when the group is not virtually free, where virtually free means having a free subgroup of finite index. This reported criterion is retained with its source scope: the excerpt does not specify finite generation, torsion, or other ambient hypotheses needed for the stated known case. Those conventions must be checked before using it as an unrestricted theorem about every discrete subgroup.

The supplied research summary credits Kahn--Markovic with closed quasi-Fuchsian surface subgroups for closed hyperbolic \(3\)-manifolds and mentions results under hypotheses for cusped manifolds, without stating those hypotheses or references. It says a full characterization in the parabolic case is not established; this is an editorial literature assessment, not an independent audit.
\subsection{Short English statement}
Characterize Kleinian groups containing a quasi-Fuchsian subgroup isomorphic to a closed genus-at-least-two surface group. The source singles out parabolic groups as difficult and reports a no-parabolic criterion, whose exact ambient hypotheses are omitted.
\subsection{Sources}
\begin{itemize}
\item Local UnsolvedMath registry, permanent ID 10900015, source code catalogue entry; attribution Agol, subsection 3.7, ``Kleinian groups with closed quasi-Fuchsian surface subgroups.''
\item Delp, Hoffoss and Manning, \emph{Problems In Groups, Geometry, and Three-Manifolds} (2015), subsection 3.7, PDF page 2. Supplied URLs: \url{https://arxiv.org/pdf/1512.04620#page=2} and \url{https://arxiv.org/abs/1512.04620}; recorded access 2026-07-29; extraction source-tex checked against rendered PDF text. Background reports the no-parabolic, not-virtually-free criterion.
\item The research summary mentions Kahn--Markovic and unspecified cusped-manifold results without full bibliographic entries or exact hypotheses.
\item Supplied status partially\_solved and PARTIAL-PROGRESS. The source presents the item as open in 2015; no complete modern primary-literature status audit is recorded. Status evidence and redistribution terms for the publicly accessible author-submitted TeX and PDF remain NEEDS\_REVIEW.
\end{itemize}
\end{document}
